\chapter{اللقاء الرابع والعشرون: آيات الإنفاق والحج والإحصار والهدي}

\section{مقدمة}
السلام عليكم ورحمة الله وبركاته. الحمد لله رب العالمين، وأشهد أن لا إله إلا الله وحده لا شريك له، وأشهد أن محمداً عبده ورسوله. اللهم صل على محمد وعلى آل محمد كما صليت على آل إبراهيم إنك حميد مجيد، اللهم بارك على محمد وعلى آل محمد كما باركت على آل إبراهيم إنك حميد مجيد.

استأنف الشيخ هذا اللقاء من قوله تعالى: \begin{ayah}وَأَنْفِقُوا فِي سَبِيلِ اللَّهِ وَلَا تُلْقُوا بِأَيْدِيكُمْ إِلَى التَّهْلُكَةِ وَأَحْسِنُوا إِنَّ اللَّهَ يُحِبُّ الْمُحْسِنِينَ\end{ayah} ثم انتقل إلى أوائل آيات الحج والإحصار والفدية في سورة البقرة، مع عناية ظاهرة بإبراز طريقة \rawi{الطبري} في تصنيف الأقوال المتقاربة، وبيان عللها، ثم ترجيح ما تؤيده الآثار واللغة والسياق.

\separator

\section{تأويل (وأنفقوا في سبيل الله ولا تلقوا بأيديكم إلى التهلكة)}
\isnad{قال أبو جعفر:} اختلف أهل التأويل في معنى التهلكة في هذه الآية على أقوال متقاربة من جهة الأصل، لكنها تختلف في جهة التعلق:
\begin{itemize}
    \item فطائفة جعلت التهلكة في ترك النفقة في سبيل الله والإمساك عنها خوف الفقر أو سوء الظن بالخلف.
    \item وطائفة جعلتها في الخروج إلى الجهاد بغير عدة ولا نفقة، فيلقي المرء نفسه إلى الهلاك بعجزه وضعفه.
    \item وطائفة حملتها على القنوط من رحمة الله بعد الذنب، فيستسلم العبد لمعصيته ويقول: لا توبة لي.
\end{itemize}

ونبّه الشيخ هنا إلى دقة \rawi{الطبري} في تصنيف الأقوال؛ فإنه لا يذيب الأقوال المتقاربة في عنوان واحد، بل يفرق بين ما اتفق في الأصل واختلف في جهة التطبيق، حتى يظهر محل النزاع بدقة.

\begin{fawaid}[تحرير الفروق بين الأقوال المتشابهة من أهم مفاتيح الفهم]
قد يظن القارئ أن الأقوال في تفسير آية ما متماثلة، بينما بينها فروق دقيقة تؤثر في المعنى والترجيح. ومن قوة صنيع الطبري أنه يميز بين هذه الفروق، فلا يخلط بين من جعل التهلكة في ترك النفقة، ومن جعلها في الخروج بلا نفقة، ومن جعلها في اليأس بعد الذنب.
\end{fawaid}

\begin{fawaid}[القنوط من رحمة الله من صور التهلكة]
من المعاني التربوية الجليلة في هذا الموضع أن العبد إذا أذنب فلا يجوز له أن يستسلم لذنبه أو يظن أن باب التوبة قد أغلق دونه؛ فإن هذا الاستسلام نفسه لون من ألوان الإلقاء بالنفس إلى التهلكة.
\end{fawaid}

\section{تأويل (وأحسنوا إن الله يحب المحسنين)}
جاء هذا الأمر عقب النفقة والجهاد تنبيهاً على أن الإحسان مطلوب في أصل العمل وفي كيفيته: إحسان في البذل، وإحسان في القصد، وإحسان في الصبر والثبات، وإحسان في رجاء الخلف من الله.

\separator

\section{تأويل (وأتموا الحج والعمرة لله)}
\isnad{قال أبو جعفر:} أمر الله عباده بإتمام الحج والعمرة له خالصين، على وجه الإكمال والإحسان والقيام بما يلزمهما من الأعمال والحدود.

ثم توسع \rawi{الطبري} في هذا الموضع توسعاً كبيراً، فذكر أقوال السلف والفقهاء في معنى الإتمام، وفي حكم العمرة: هل هي فرض لازم كالحج أم هي تطوع مؤكد؟ وساق لذلك آثاراً وأخباراً، ونظر في ثبوتها، ووازن بين ظواهر الأدلة.

والذي يظهر من صنيعه هنا أن الكلام ليس مجرد تفريع فقهي مستقل، بل هو تفسير للآية لا يتم إلا بتحرير هذه المسائل؛ لأن فهم الأمر بالإتمام متوقف على معرفة ما الذي دخل فيه المكلف ابتداءً، وما الذي يجب عليه إتمامه، وما الفرق بين أصل الوجوب وبين وجوب الإتمام بعد الشروع.

\begin{fawaid}[ليس كل ما يذكره الطبري في الفقه خروجاً عن التفسير]
إذا توسع الطبري في مسألة فقهية أو لغوية فالغالب أن ذلك لأنه يرى أن فهم الآية لا يتم إلا بها. فذكره للخلاف هنا في الحج والعمرة والإحصار والهدي ليس استطراداً مجرداً، بل هو من تمام بيان مدلول الآية.
\end{fawaid}

\section{الخلاف في وجوب العمرة}
عرض \rawi{الطبري} في هذه الآية جملة من الآثار في وجوب العمرة، ونظر في ثبوت المرفوع منها والموقوف، مع مقارنتها بظاهر الآية وبالآثار الأخرى. واعتنى الشيخ هنا بالتنبيه إلى أن من أهم ما يميز صنيع \rawi{الطبري} أنه لا يبني تقرير التفسير على مجرد خبر لا يثبت، بل يذكر ضعفه أو يعلّق المعنى على صحته.

ومن خلال مجموع كلامه يظهر ميله إلى أن العمرة ليست في رتبة الحج من جهة الفرضية الأصلية، وإن كان الشروع فيها يوجب إتمامها، وإن كان فضلها ومشروعيتها مؤكدين.

\begin{fawaid}[صحة الخبر شرط في بناء التفسير عليه]
من أهم معالم منهج الطبري أنه لا يسوي بين الأخبار في مقام الاستدلال. فقد يذكر الحديث ثم يبين وهاءه أو يعلّق الاستدلال به على ثبوته، وبذلك يحفظ باب التفسير من أن يبنى على أخبار لا تقوم بها حجة.
\end{fawaid}

\separator

\section{تأويل (فإن أحصرتم فما استيسر من الهدي)}
\isnad{قال أبو جعفر:} إذا منع المحرم من إتمام نسكه بحج أو عمرة، فعليه ما تيسر من الهدي ليحل من إحرامه.

وهنا توسع \rawi{الطبري} في تحرير معنى الإحصار: هل يختص بعدو يمنع من الوصول إلى البيت؟ أم يدخل فيه المرض، والكسر، والخوف، وضلال الطريق، وكل ما يحبس المحرم عن إتمام نسكه؟

ونقل في ذلك أقوالاً كثيرة عن \rawi{مالك} و\rawi{ابن عمر} وغيرهما، وظهر من مجموع عرضه أنه يميز بين مواضع الخلاف تمييزاً دقيقاً:
\begin{itemize}
    \item خلاف في معنى الإحصار نفسه.
    \item وخلاف في محل الهدي: أهو موضع الحبس، أم لا يبلغ محله حتى يصل إلى الحرم؟
    \item وخلاف في الفرق بين الحج والعمرة في باب الإحصار.
    \item وخلاف في القضاء بعد التحلل: هل يجب مطلقاً أو في بعض الصور؟
\end{itemize}

ثم رجح \rawi{الطبري} أن الإحصار يعم كل محصر في حج أو عمرة، وأن محل هديه الموضع الذي أحصر فيه، حلاً كان أو حرماً، وأن له أن يتحلل إذا بلغ الهدي محله، مستنداً في ذلك إلى الأخبار المتواترة في قصة الحديبية، وإلى كون الآية نزلت في هذا السياق.

\begin{fawaid}[التنزيل أولاً ثم التنزل]
منهج الطبري المحكم أن يبدأ بسبب النزول وما نزلت فيه الآية ابتداءً، ثم يمد الحكم إلى ما يشبهه بدليل صحيح. ولذلك جعل قصة الحديبية أصلاً في باب الإحصار، ثم نظر في سائر الصور على ضوء هذا الأصل.
\end{fawaid}

\begin{fawaid}[فعل النبي صلى الله عليه وسلم من أقوى ما يفسر به القرآن]
لما كانت آية الإحصار قد نزلت في سياق صلح الحديبية، وكان فعل النبي صلى الله عليه وسلم فيها ثابتاً مشهوراً، جعل الطبري هذا الفعل أرجح ما تفسر به الآية، وقدم العمل النبوي الثابت على الأقوال المتجردة عنه.
\end{fawaid}

\section{تأويل (فما استيسر من الهدي)}
تكلم \rawi{الطبري} على أقل ما يجزئ في هذا الهدي، ونظر في معنى \textit{ما استيسر}، وهل المراد به مطلق ما تيسر من النعم، أو أن أقل ذلك شاة، وما الذي يثبت بالدليل من هذه الصور.

والذي يظهر من صنيعه أنه يميل إلى الاكتفاء بما يصدق عليه اسم الهدي مما جاءت به دلالة الشرع، مع بناء الترجيح على ظاهر التنزيل وما صح من البيان، لا على مجرد التوسع في الاحتمالات.

\section{تأويل (ولا تحلقوا رؤوسكم حتى يبلغ الهدي محله)}
بيّن \rawi{الطبري} أن حلق الرأس علامة التحلل من الإحرام، فلا يجوز للمحصر أن يتحلل حتى يبلغ الهدي محله، ثم وقع الخلاف في تفسير هذا المحل:
\begin{itemize}
    \item فقيل: محل الهدي موضع الحبس نفسه إذا ذبح أو نحر فيه، خاصة إذا كان الإحصار بعدو.
    \item وقيل: لا يبلغ الهدي محله حتى يصل إلى الحرم.
    \item وقيل بتفصيلات أخرى بحسب نوع الإحصار وحال النسك.
\end{itemize}

ثم رجح \rawi{الطبري} أن المحل هو الموضع الذي يحل فيه ذبحه أو نحره، وأن ذلك قد يكون في موضع الإحصار نفسه، مستدلاً بقصة الحديبية وما وقع فيها من نحر الهدي ثم الحلق والتحلل.

\begin{fawaid}[من قوة الاستدلال جمع الآية إلى السنة إلى السيرة]
الترجيح عند الطبري لا يبنى على مجرد احتمال لغوي، بل يجمع بين دلالة اللفظ القرآني، والسنة الثابتة، والسياق التاريخي الذي نزلت فيه الآية. وهذه الثلاثة اجتمعت هنا في باب الإحصار والهدي.
\end{fawaid}

\separator

\section{تأويل (فمن كان منكم مريضاً أو به أذى من رأسه ففدية من صيام أو صدقة أو نسك)}
وقف الشيخ عند أول هذه الآية، وأشار إلى أن \rawi{الطبري} ابتدأ فيها بذكر الخلاف في معنى الفدية وأنواعها، ومقدار الصيام والصدقة والنسك، إلا أن التفصيل المؤجل في هذا اللقاء كان مركزاً على ما سبقها من آيات الإحصار والهدي.

ومع ذلك نبه في خاتمة الدرس إلى أهمية جمع مواضع الخلاف في هذا المقطع كله من قوله تعالى \begin{ayah}وَأَتِمُّوا الْحَجَّ وَالْعُمْرَةَ لِلَّهِ\end{ayah} إلى أول الكلام على الفدية؛ لأن هذا الموضع من أوسع المواضع التي يظهر فيها منهج \rawi{الطبري} في:
\begin{itemize}
    \item استيعاب الأقوال.
    \item تحرير صورة النزاع.
    \item بيان علة كل مذهب.
    \item ترجيح ما يسنده الأثر الصحيح واللسان والسياق.
\end{itemize}

\begin{fawaid}[من يقرأ الطبري ليتتبع طريقته ينتفع أعظم من مجرد قارئ النتائج]
الفرق كبير بين من يقرأ ليعرف القول المختار فقط، وبين من يقرأ ليفهم كيف نشأ الخلاف، وكيف صُنّفت الأقوال، وكيف رُجِّح بينها. ومن هنا كان هذا الموضع من سورة البقرة من أنفس المواضع لاكتشاف المنهج العملي للإمام الطبري.
\end{fawaid}

\separator

ختم الشيخ اللقاء بالتأكيد على أن هذا المقطع من أخصب مواضع الكتاب في التدريب على قراءة تفسير \rawi{الطبري} قراءة منهجية، لا تكتفي بحفظ النتائج، بل ترصد طريقته في جمع الأقوال، وتحليلها، وتخريجها، والاعتراض عليها، ثم اختيار ما يراه أقرب إلى الدليل.
